% ch00  引言（原文 555–726 行）

\section*{引言}
\addcontentsline{toc}{section}{引言}

\subsection{一些动机}

设 $F$ 是完备离散赋值域，$\mO\subset F$ 是其整数环，$k$ 是剩余域。又设 $V=F^n$ 是 $n$ 维 $F$-向量空间。
$V$ 的一个\emph{格}（lattice）是指 $V$ 的有限生成 $\mO$-子模 $\La$，且满足 $\La\otimes F=V$。例如 $\La_0=\mO^n$ 就是 $V$ 中的一个格。
本讲义的主要目的，是给 $V$ 的格构成的集合（及其变形与推广）赋予无穷维代数簇的结构，即所谓的\emph{仿射 Grassmann 簇} $\Gr$；随后研究它的几何性质，并给出若干应用。

不过，我们先给出研究这一对象的一份不完整的动机清单。

第一个动机来自环路代数、环路群与 Kac--Moody 代数、Kac--Moody 群之间的关系。由于每个格都可以经 $V$ 的线性自同构搬到 $\La_0$，
作为集合，$\Gr$ 自然等同于 $\GL_n(F)/\GL_n(\mO)$。当 $F=k((t))$ 时，群 $\GL_n(F)$ 与 $k$ 上的仿射 Kac--Moody 群密切相关。
把 $\Gr$ 实现为该 Kac--Moody 群的（部分）旗簇，对仿射 Kac--Moody 代数的表示论很重要。两者的确切关系将在 \S\ref{Kac-Moody} 中简述。

第二个动机来自代数曲线上向量丛的模空间。设 $X$ 是 $k$ 上光滑投影代数曲线，$x\in X(k)$ 是一点。取 $x$ 附近的局部坐标 $t$，
把 $X$ 在 $x$ 处的完备局部环与 $k[[t]]$ 等同。于是可以把 $\Gr$ 等同于下述集合
\[
\Gr=\left\{\, (\mE,\beta) \ \middle|\ \mE \text{ 是 } X \text{ 上的向量丛},\ \beta:\mE|_{X-x}\simeq \mO_{X-x}^n \text{ 是平凡化} \,\right\}.
\]
于是有从 $\Gr$ 到 $X$ 上秩 $n$ 向量丛的同构类集合的映射。后者实际上是一个精细得多的代数几何对象的点集，即 $X$ 上秩 $n$ 向量丛的模叠 $\Bun_n$。
因此很自然希望把上述集合论层面的映射升级为代数几何对象之间的态射
\[
u_x:\Gr\to \Bun_n .
\]
这样的态射称为\emph{单值化映射}（uniformization map），它是研究 $\Bun_n$ 的某些几何性质的重要工具。关于这一想法的详细解释，见 \S\ref{one point uniformization}。

第三个动机来自 $\Gr$ 中一些（局部）闭子簇的研究，它们出现于朗兰兹纲领的各个方面。这类子簇包括仿射 Springer 纤维（与非 Archimedes 局部域上约化群的调和分析有关）、
仿射 Deligne--Lusztig 簇（与志村簇模 $p$ 的约化有关）以及 Kisin 簇（与 Galois 表示的形变空间有关）。这些簇的简介见 \S\ref{digression}。

第四个动机同样来自朗兰兹纲领。描述 $\GL_n(F)$ 上 $\GL_n(\mO)$-双不变函数环的 Satake 同构\zy{更确切地说，是 Langlands 对 Satake 同构的重新解释。}，
是朗兰兹对偶的出发点。事实证明，Satake 同构有一个极大的推广，称为\emph{几何 Satake 等价}。它是几何朗兰兹纲领的出发点，
近年来还在经典算术朗兰兹纲领中找到了一些出人意料的用途（例如见 \cite{Laf}）。但即便只是为了陈述几何 Satake 等价，我们也需要知道 $\Gr$ 上的代数几何结构。

\subsection{范围与内容}

\subsubsection{} 本讲义以 2015 年夏在 Park City 数学研究所的五次讲座为基础。它的定位是对本课题的一次严肃导引。
话虽如此，我们只处理（我们认为）本课题最基本也最重要的方面。本讲义当然没有触及许多其他重要课题，一部分是因为篇幅有限，一部分是因为作者能力不足。

另一方面，我们希望本讲义是对既有文献的补充，而不是取代它们。因此，对每个结果我们并不都给出详细证明，尤其是当文献中已有这样的证明时；
这种情况下我们只指出相关文献。反过来，有些众所周知且被广泛使用的结果，文献中似乎并没有（完整的）证明。遇到这种情况，我们会尽量补出细节。

本讲义包含少数此前未见于文献的结果，尽管它们很可能已为专家所熟知。例如 Beilinson--Drinfeld Grassmann 簇的相对 Picard 层的计算（\S\ref{line bundle on BD Grass}）。

\subsubsection{} 下面简述各节内容。

第 \ref{Lecture I} 节解释基本构造，包括仿射 Grassmann 簇的各种（模）解释。我们还将引入重要的行列式线丛，
并讨论复数域上的仿射 Grassmann 簇与基于环路群之间的关系。虽然我们主要在等特征的局部域上考虑群，
但在本节末尾会简述为 $p$-进群定义仿射 Grassmann 簇的近期进展。

第 \ref{Lecture II} 节研究仿射 Grassmann 簇的一些基本几何。我们将引入仿射 Grassmann 簇中的各种子簇，包括 Schubert 簇、对偶 Schubert「簇」、横截片，
以及一些源于朗兰兹纲领的重要子簇。对单连通单群 $G$，我们研究其仿射 Grassmann 簇的 Picard 群。
我们还简述仿射 Grassmann 簇与仿射 Kac--Moody 群的部分旗簇之间的关系。

把仿射 Grassmann 簇解释为曲线上丛的某种模空间，这一想法有一个极大的推广，即 Beilinson--Drinfeld Grassmann 簇。
第 \ref{Lecture III} 节专门讨论这些对象。Beilinson--Drinfeld Grassmann 簇的一个显著特征是\emph{分解性质}（factorisation property）。
我们较详细地解释这一性质，特别是用曲线的 Ran 空间来表述它。我们还讨论 Beilinson--Drinfeld Grassmann 簇上的线丛。

第 \ref{Lecture IV} 节讨论仿射 Grassmann 簇对曲线上丛模空间研究的一些应用。基本工具是单值化定理。这样的定理实际上有两类：
当 $G$ 半单时成立的单点单值化，以及 ad\'elique 单值化——它是 Weil 对曲线上 $G$-丛同构类的 ad\`ele 描述的代数几何版本。
作为应用，当 $G$ 单且单连通时，我们给出光滑投影曲线上 $G$-丛模空间的 Picard 群的描述，并把模空间上线丛的整体截面解释为共形块（conformal block）。

第 \ref{Lecture V} 节专门讨论几何 Satake 等价。这一理论自然分成两部分。第一部分是从仿射 Grassmann 簇上某类反常层构造一个 Tannaka 范畴。
关键要素是该范畴上的\emph{融合积}结构，它的构造依赖 Beilinson--Drinfeld Grassmann 簇的分解性质。
第二部分是研究 Tannaka 群。这里的主要工具是 Mirkovi\'c--Vilonen 关于半无限轨道几何的理论。
我们还解释如何把几何 Satake 等价推广到适合（几何）朗兰兹纲领应用的形式。本节末尾，我们说明如何从几何版本恢复经典 Satake 同构。

另外有一个附录，收集第 \ref{Lecture V} 节用到的 $\ell$-进层的若干背景知识。

\subsection{约定与记号}

下面固定全讲义通用的记号与约定。读者\emph{务必}先跳过本小节，只在需要时回来查阅。

\subsubsection{$k$-空间}\label{Intro:sp}

设 $k$ 是交换域。以 $\on{Aff}_k$ 记仿射 $k$-概形范畴，即交换 $k$-代数范畴 $k\on{-alg}$ 的反范畴。
\emph{预层}（prestack 叫\emph{预叠}）是指从 $\Aff_k$ 到集合范畴（相应地，群胚 $2$-范畴）的共变函子（相应地，$2$-函子）。
我们把集合看成一个离散群胚，因此有时把预层看成离散的预叠。

给 $\Aff_k$ 赋以 fpqc 拓扑，即 $R$ 的一个覆盖是忠实平坦映射 $R\to R'$。$k$-\emph{空间}是指关于该拓扑在 $\Aff_k$ 上的层。
具体地说，空间 $\mF$ 是共变函子 $k\on{-alg}\to \on{Set}$，满足
\begin{enumerate}
\item 保持有限乘积，即
\[
\mF\Big(\prod_i R_i\Big)\to \prod_i \mF(R_i)
\]
是同构；
\item 且当 $R\to R'$ 忠实平坦时，
\begin{equation}\label{fpqc sheaf}
\mF(R)\to \mF(R')\rightrightarrows\mF(R'\otimes_RR')
\end{equation}
是等化子。
\end{enumerate}
两个空间之间的态射即函子的自然变换。$k$-空间范畴记为 $\Sp_k$。众所周知，$\Sp_k$ 以 $k$-概形范畴 $\on{Sch}_k$ 为全子范畴。
回忆：预层映射 $f:\mF\to\mG$ 称为\emph{概形化的}（schematic），若对每个概形 $T$，纤维积 $\mF\times_{\mG}T$ 都是概形。叠的概念可以类似地定义。

设 $X$ 与 $Y$ 是两个预叠。纤维积 $X\times Y$ 总是指 $X\times_kY$。若 $Y=\Spec R$，有时也把它写成 $X_R$、$X\otimes R$ 或 $X\otimes_kR$。
当 $Y=\Spec A$、$Z=\Spec B$ 时，我们也把 $X\times_YZ$ 写成 $X\otimes_AB$。

设 $X$ 是 $k$ 上的预叠。定义 $X$ 上的正则函数环为
\begin{equation}\label{reg fun}
\Gamma(X,\mO_X)=\Hom(X, \bG_a),
\end{equation}
其中 $\Hom$ 在预叠范畴中计算。也就是说，$X$ 上的正则函数 $f$ 是这样一种相容的方式：给每个 $x\in X(R)$ 指定元素 $f_x\in R$，并满足自然的相容性条件。
类似地，定义 $X$ 上线丛的 Picard 群胚为
\begin{equation}\label{Pic line}
\Pic(X)=\Hom(X,\bB\bG_m),
\end{equation}
其中 $\bB\bG_m$ 是 $\bG_m$ 的分类叠（见下面 \S\ref{Intro:torsor}）。也就是说，$\Pic(X)$ 的对象（称为 $X$ 上的线丛）是这样一种相容的方式：
给每个 $x\in X(R)$ 指定 $\Spec R$ 上的一个线丛，并满足自然的相容性条件。

\subsubsection{归纳概形}\label{Intro:Indsch}

本讲义中有一类特殊的 $k$-空间扮演重要角色。

\begin{dfn}\label{ind-sch}
$k$ 上的\emph{归纳概形}（ind-scheme）是指可写成子空间的\emph{滤过}正向极限
\begin{equation}\label{presentation}
X=``\varinjlim''_{i\in I}X_i
\end{equation}
的空间，其中每个 $X_i$ 都可由概形表示，且对 $I$ 中每个箭头 $i\to j$，相应映射 $X_i\to X_j$ 都是闭浸入。
\end{dfn}

\begin{rmk}
文献中，上面定义的归纳概形有时称为\emph{严格}归纳概形，以区别于概形范畴中一般的归纳对象。后一概念在本讲义中不会出现。
\end{rmk}

\begin{dfn}
设 $X$ 是 $k$ 上的归纳概形。称 $X$ 为：(i) \emph{归纳仿射的}；(ii) \emph{在 $k$ 上归纳有限型的}；(iii) \emph{在 $k$ 上归纳 proper 的}；
(iv) \emph{在 $k$ 上归纳投影的}；(v) \emph{既约的}，若存在形如 \eqref{presentation} 的表示，使得每个 $X_i$ 都可取为
(i) 仿射的；(ii) 在 $k$ 上有限型的；(iii) 在 $k$ 上 proper 的；(iv) 在 $k$ 上投影的；(v) 既约的。
定义 $X_{\on{red}}:=``\varinjlim''_{i\in I}(X_i)_{\on{red}}$。
\end{dfn}

\begin{ex}
形式概形是这样一种归纳概形：$X_{\on{red}}$ 是概形。
\end{ex}

\subsubsection{主丛}\label{Intro:torsor}

设 $G$ 是 $k$-群，即 $k$-空间范畴中的群对象。预叠映射 $E\to X$ 称为 $G$-\emph{主丛}（$G$-torsor），若 $G$ 作用在 $E$ 上，并且：
\begin{enumerate}
\item 映射 $E\to X$ 是 $G$-等变的（其中 $G$ 平凡地作用在 $X$ 上）；
\item 对每个 $\Spec R\to X$，存在忠实平坦映射 $\Spec R'\to\Spec R$，使得 $\Spec R'\times_XE$ 与 $\Spec R'\times G$ 是 $G$-等变同构的。
\end{enumerate}
若 $G$ 是 $k$ 上仿射群概形，这等价于要求 $E\to X$ 概形化且忠实平坦，且自然作用映射 $G\times E\cong E\times_X E$ 是同构。
若 (2) 中的 $\Spec R'$ 可取成使 $\Spec R'\to\Spec R$ 为 fppf（相应地，\'etale、Zariski）覆盖，则称该 $G$-主丛是 fppf（相应地，\'etale、Zariski）局部平凡的。

设 $X$ 是带 $k$-群 $G$ 作用的 $k$-空间。以 $[X/G]$ 记 fpqc 叠，其 $R$-点由 $\Spec R$ 上的 $G$-主丛 $E$ 连同 $G$-等变映射 $E\to X$ 组成。
自然投影 $X\to[X/G]$ 便是 $G$-主丛。若 $X=\Spec k$，通常把 $[X/G]$ 记为 $\bB G$，称为 $G$ 的分类叠。若 $G$ 在 $X$ 上的作用自由，则 $[X/G]$ 是 $k$-空间。
设 $E\to X$ 是 $G$-主丛，$Y$ 是 $G$-空间（即带 $G$ 作用的 $k$-空间），定义\emph{扭曲积}为
\[
X\tilde\times Y:= E\times^G Y:=[E\times Y/G],
\]
其中 $G$ 对角地作用在 $E\times Y$ 上。有自然投影 $X\tilde\times Y\to X$，其纤维同构于 $Y$。若 $Y=V$ 是 $G$ 的线性表示，
有时也把 $E\times^GV$ 写成 $V_E$，并把它看成 $X$ 上的向量丛。

若 $G=\GL_n$ 是一般线性群，$V=k^n$ 是其标准表示，则 $(\mE\to X)\mapsto(V_\mE\to X)$ 给出 $X$ 上 $\GL_n$-主丛与 $X$ 上秩 $n$ 向量丛的群胚等价。
本讲义中将自由使用这一等价。

上述记号显然可以推广到某个基 $k$-概形 $S$ 上的（非常值）群空间 $\underline G$。以 $\mE^0$ 记\emph{那个}平凡 $\underline G$-主丛。

\subsubsection{圆盘}\label{Intro:disc}

记
\[
D=\Spec k[[t]],\quad D^*=\Spec k((t)),
\]
称之为\emph{那个}\emph{圆盘}与\emph{去心圆盘}，以强调坐标 $t$ 的选定。对 $k$-代数 $R$，记
\[
D_R:= D\hat\times\Spec R:=\Spec R[[t]], \quad D_R^*:= D^*\hat\times \Spec R:= \Spec R((t)),
\]
并把它们看成由 $\Spec R$ 参数化的圆盘族（相应地，去心圆盘族）。

除另有说明，所谓局部域 $F$，指 $k$ 上的域，它关于一个在 $k$ 上平凡的离散赋值完备，且剩余域为 $k$。以 $\mO\subset F$ 记其整数环。
于是选定素元 $t\in\mO$ 后，有 $F\cong k((t))$、$\mO\cong k[[t]]$。设
\[
\mO\hat\otimes R:=\varprojlim (\mO/t^n\otimes R), \quad F\hat\otimes R:=\mO\hat\otimes R[1/t],
\]
从而 $\mO\hat\otimes R\cong R[[t]]$、$F\hat\otimes R\cong R((t))$。记 $D_{F,R}=\Spec (\mO\hat\otimes R)$，$D^*_{F,R}=\Spec (F\hat\otimes R)$。

设 $X$ 是 $k$ 上光滑几何连通曲线。以 $|X|$ 记 $X$ 的闭点集。对 $x\in|X|$，记 $X^*=X-\{x\}$。
以 $\mO_x$ 记 $X$ 在 $x$ 处的完备局部环，$F_x$ 记 $\mO_x$ 的分式域。对 $k$-代数 $R$，为简单起见把 $D_{F_x,R}$、$D^*_{F_x,R}$ 记成 $D_{x,R}$、$D^*_{x,R}$。
若 $x:\Spec R\to X$ 是 $X$ 的 $R$-点，记 $\Ga_x\subset X_R$ 为 $x$ 的图像，$\hat{\Ga}_x$ 为 $\Ga_x$ 在 $X_R$ 中的形式完备化。

\subsubsection{群论数据}\label{Intro:group}

设 $G$ 是 $k$ 上连通分裂约化群。以 $G_\der$ 记 $G$ 的导群，$G_\ad$ 记 $G$ 的伴随形式，$G_\s$ 记 $G_\der$ 的单连通覆盖。

以 $T$ 记 $G$ 的抽象 Cartan。回忆它定义为 Borel 子群 $B\subset G$ 对其幂幺根群的商。可以证明 $T$ 在不计典范同构的意义下与 $B$ 的选取无关。
当需要把 $T$ 作为极大环面嵌入 $G$ 时，我们会明确说明。

以 $\xch=\xch(T)$ 记权格，即自由 Abel 群 $\Hom(T,\bG_m)$；以 $\xcoch=\xcoch(T)$ 记余权格，即 $\xch(T)$ 的对偶。
以 $\Phi\subset\xch$（相应地，$\Phi^\vee\subset\xcoch$）记根系（相应地，余根系）。四元组 $(\xch,\xcoch,\Phi,\Phi^\vee)$ 称为 $G$ 的根数据。
以 $W$ 记 Weyl 群。Borel 子群 $B\subset G$ 决定了支配余权构成的半群 $\xcoch(T)^+\subset\xcoch(T)$ 与正根子集 $\Phi^+\subset\Phi$，它们其实与 $B$ 的选取无关。于是有典范双射
\[
\xcoch(T)/W\cong \xcoch(T)^+.
\]
回忆 $\xcoch(T)$ 上有偏序 $\leq$：若 $\mu-\la$ 是单余根的非负整系数线性组合，则记 $\la\leq\mu$；若 $\la\leq\mu$ 但 $\la\neq\mu$，则记 $\la<\mu$。
该偏序在 $\xcoch(T)^+$ 上的限制有时称为 $\xcoch(T)^+$ 上的 Bruhat 序。通常以 $2\rho\in\xch(T)$ 记所有正根之和，并令 $\rho=\frac12(2\rho)\in\xch(T)\otimes\bQ$。

\subsection{致谢}

本讲义基于 2015 年夏在 Park City 数学研究所的讲座。我要感谢 PCMI 的组织者、听众与工作人员。
我还要感谢Weizhe Zheng就 \S\ref{ULA sheaf} 的有益讨论，以及 Michael Harris、Jiuzu Hong与 Joel Kamnitzer 的宝贵意见。
